\documentclass[11pt,a4paper]{article}
\usepackage[margin=25mm,headheight=14pt]{geometry}
\usepackage[T1]{fontenc}
\usepackage[utf8]{inputenc}
\usepackage{lmodern,amsmath,amssymb}
\usepackage{xcolor,graphicx,booktabs,longtable,array}
\usepackage{microtype,enumitem,fancyhdr,titlesec}
\usepackage{hyperref}
\definecolor{accent}{HTML}{163B55}
\hypersetup{colorlinks=true,linkcolor=accent,urlcolor=accent,
 pdftitle={Critical inradius asymptotics and a supremum-norm bound for convex domains},pdfauthor={Matthew J. Colbrook},
 pdfsubject={AIM 1016: mathematical result and proof}}
\urlstyle{same}
\Urlmuskip=0mu plus 1mu
\setlength{\parindent}{0pt}
\setlength{\parskip}{0.35em}
\setlength{\emergencystretch}{3em}
\setlist{itemsep=0.25em,topsep=0.3em}
\providecommand{\tightlist}{\setlength{\itemsep}{0pt}\setlength{\parskip}{0pt}}
\titleformat{\section}{\normalfont\large\bfseries\color{accent}}{\thesection}{0.7em}{}
\titleformat{\subsection}{\normalfont\normalsize\bfseries}{\thesubsection}{0.6em}{}
\pagestyle{fancy}
\fancyhf{}
\fancyhead[L]{\small\sffamily AIM 1016}
\fancyhead[R]{\small\sffamily Matthew J. Colbrook}
\fancyfoot[C]{\thepage}
\renewcommand{\headrulewidth}{0.3pt}
\setcounter{secnumdepth}{0}
\allowdisplaybreaks[2]
\title{\textbf{Critical inradius asymptotics and a supremum-norm bound for convex domains}}
\author{Matthew J. Colbrook\\[3pt]
{\small Department of Applied Mathematics and Theoretical Physics}\\
{\small University of Cambridge, United Kingdom}\\
{\small\href{mailto:mjc249@maths.cam.ac.uk}{\texttt{mjc249@maths.cam.ac.uk}}}}
\date{3 November 2026}
\begin{document}
\maketitle
\begin{abstract}
We determine the leading large-exponent asymptotic of the critical inradius-normalized Dirichlet frequency constant. We also give an explicit lower bound for the supremum-norm problem in a separate range of exponents. These results address parts of AIM 1016 and do not determine its general sharp constants.
\end{abstract}
\noindent\textbf{Repository target.}
This manuscript concerns AIM problem 1016 at revision
\texttt{3fa50884d129}. Its scope is stated below.
\medskip

\hypertarget{1016-result-and-scope}{%
\section{Result and scope}\label{1016-result-and-scope}}

For a bounded nonempty open convex domain \(\Omega\subset\mathbb R^N\),
let \(R_\Omega\) be its Euclidean inradius. In the critical case
considered here, the constant in AIM 1016 is

\[
C_{\mathrm{opt}}(N,N,q)=
\inf_{\Omega\ \mathrm{convex}} R_\Omega^{N^2/q}
\inf_{\substack{u\in C_c^\infty(\Omega)\\u\ne0}}
\frac{\int_\Omega|\nabla u|^N\,dx}
 {(\int_\Omega|u|^q\,dx)^{N/q}}.
\]

Equivalently, impose \(R_\Omega=1\) and take the infimum of the
displayed quotient. Write

\[
\sigma_{N-1}=|S^{N-1}|,\qquad N'=\frac{N}{N-1}.
\]

For every integer \(N\ge 2\), the following asymptotic holds:

\[
\boxed{\displaystyle
\lim_{q\to\infty}q^{N-1}C_{\rm opt}(N,N,q)
=\sigma_{N-1}\left(\frac{eN^2}{N-1}\right)^{N-1}.}
\tag{1}
\]

In particular,

\[
\lim_{q\to\infty}q\,C_{\rm opt}(2,2,q)=8\pi e.
\]

This determines an asymptotic regime of the parameter-dependent problem,
\textbf{not} the exact constant at a fixed finite \(q\). The proof is a
consequence of classical sharp exponential integrability, a uniform
convex-domain Hardy estimate, and an explicit concentrating test. No
claim that the asymptotic is new in the literature is made.

\hypertarget{1016-a-shape-uniform-ln-estimate}{%
\section{1. A shape-uniform L\^{}N
estimate}\label{1016-a-shape-uniform-ln-estimate}}

By scaling, restrict to domains of inradius one. Their boundary-distance
function \(d(x)=\operatorname{dist}(x,\partial\Omega)\) satisfies
\(0<d\le 1\), \(|\nabla d|=1\) almost everywhere, and is concave because
the domain is convex. In particular \(\Delta d\le 0\) distributionally.

For \(u\in C_c^\infty(\Omega)\), integration by parts using

\[
\operatorname{div}\left(\frac{\nabla d}{d^{N-1}}\right)
\le-\frac{N-1}{d^N}
\]

and then Hölder's inequality give the standard convex Hardy estimate

\[
\left\|\frac{u}{d}\right\|_N
\le \frac{N}{N-1}\|\nabla u\|_N.
\]

The distributional computation can equivalently use the locally BV
field \(\nabla d/(d+\varepsilon)^{N-1}\), whose singular divergence
part is nonpositive, and then let \(\varepsilon\downarrow0\).
Adding \(\varepsilon\) regularizes the denominator, not the distance
function itself. Since \(d\le 1\), this implies

\[
\|u\|_N^N\le B_N\|\nabla u\|_N^N,
\qquad B_N=\left(\frac{N}{N-1}\right)^N.
\tag{2}
\]

The constant is uniform over the entire inradius-one class; no bound on
volume or diameter is needed.

\hypertarget{1016-exponential-integrability-and-moment-bounds}{%
\section{2. Exponential integrability and moment
bounds}\label{1016-exponential-integrability-and-moment-bounds}}

Use the Adachi--Tanaka inequality on \(\mathbb R^N\). Put

\[
\alpha_N=N\sigma_{N-1}^{1/(N-1)},\qquad
\Phi_N(t)=e^t-\sum_{j=0}^{N-2}\frac{t^j}{j!}.
\]

For each fixed \(0<\alpha <\alpha _N\), there is a finite
\(A_{N,\alpha}\) such that every \(u\in W^{1,N}(\mathbb R^N)\) with
\(\Vert \nabla u\Vert _N=1\) satisfies

\[
\int_{\mathbb R^N}\Phi_N(\alpha|u|^{N'})\,dx
\le A_{N,\alpha}\|u\|_N^N.
\tag{3}
\]

This is Theorem 0.1 of the authors' 1999 primary manuscript and the
principal theorem of their 2000 journal paper:

\begin{itemize}
\tightlist
\item
  S. Adachi and K. Tanaka, \emph{A scale-invariant form of
  Trudinger--Moser inequality and its best exponent}, RIMS Kokyuroku
  \textbf{1102} (1999), 148--153. Primary PDF:
  \url{https://www.kurims.kyoto-u.ac.jp/~kyodo/kokyuroku/contents/pdf/1102-11.pdf}
\item
  S. Adachi and K. Tanaka, \emph{Trudinger type inequalities in R\^{}N
  and their best exponents}, Proc. Amer. Math. Soc. \textbf{128} (2000),
  2051--2057. DOI: \url{https://doi.org/10.1090/S0002-9939-99-05180-1}
\end{itemize}

Extend a compactly supported domain test by zero to \(\mathbb R^N\) and
normalize its gradient norm to one. Combining (2) and (3), and retaining
the \(m\)th nonnegative exponential-series term for any integer
\(m\ge N-1\), gives

\[
\int_\Omega |u|^{N'm}\,dx
\le A_{N,\alpha}B_N\,m!\,\alpha^{-m}.
\tag{4}
\]

Set \(r_m=N'm\). Taking the power \(N/r_m=(N-1)/m\) in (4), and then
taking the reciprocal, proves the shape-uniform lower bound

\[
C_{\rm opt}(N,N,r_m)
\ge \alpha^{N-1}
(A_{N,\alpha}B_N)^{-(N-1)/m}
(m!)^{-(N-1)/m}.
\tag{5}
\]

Stirling's formula yields

\[
\liminf_{m\to\infty}r_m^{N-1}C_{\rm opt}(N,N,r_m)
\ge(eN'\alpha)^{N-1}.
\tag{6}
\]

This also holds along all real \(q\to \infty\), not only the sequence
\(r_m\). Indeed, choose \(m=\lceil q/N'\rceil\), so \(q\le r_m<q+N'\),
and interpolate between the uniform \(L^N\) estimate (2) and (4):

\[
\frac1q=\frac{\theta_q}{N}+\frac{1-\theta_q}{r_m},
\qquad
\theta_q=\frac{N(r_m-q)}{q(r_m-N)}=O(q^{-2}).
\]

Thus \(\|u\|_q^N\le B_N^{\theta_q}\|u\|_{r_m}^{N(1-\theta_q)}\); the
interpolation factors are asymptotically one after division by
\(q^{N-1}\). Taking the infimum over all shapes and tests preserves this
uniform bound. Finally let \(\alpha\) increase to \(\alpha _N\) in (6),
obtaining

\[
\liminf_{q\to\infty}q^{N-1}C_{\rm opt}(N,N,q)
\ge(eN'\alpha_N)^{N-1}
=\sigma_{N-1}\left(\frac{eN^2}{N-1}\right)^{N-1}.
\tag{7}
\]

The order of limits is important: \(\alpha <\alpha _N\) is fixed while
\(q\) tends to infinity. No uniform control of \(A_{N,\alpha}\) at the
critical exponent is assumed.

\hypertarget{1016-a-matching-moser-test-in-the-unit-ball}{%
\section{3. A matching Moser test in the unit
ball}\label{1016-a-matching-moser-test-in-the-unit-ball}}

The unit ball has inradius one, so its test functions provide upper
bounds for the shape infimum. For \(L>0\), define the radial function

\[
u_L(r)=\sigma_{N-1}^{-1/N}
\begin{cases}
L^{(N-1)/N},&0\le r\le e^{-L},\\
L^{-1/N}\log(1/r),&e^{-L}<r<1,\\
0,&r\ge1.
\end{cases}
\]

This belongs to \(W_0^{1,N}(B_1)\), and direct integration gives

\[
\int_{B_1}|\nabla u_L|^N\,dx=1.
\tag{8}
\]

Its inner plateau alone yields

\[
\int_{B_1}|u_L|^q\,dx
\ge\frac{\sigma_{N-1}}{N}e^{-NL}
\sigma_{N-1}^{-q/N}L^{q(N-1)/N}.
\]

Therefore

\[
C_{\rm opt}(N,N,q)
\le \sigma_{N-1}
\left(\frac N{\sigma_{N-1}}\right)^{N/q}
e^{N^2L/q}L^{-(N-1)}.
\tag{9}
\]

The nonsmooth test is legitimate by smooth approximation in
\(W_0^{1,N}(B_1)\) and \(L^q(B_1)\) for each finite \(q\).
Alternatively, place the same test inside a fixed smaller concentric
ball and let its radius tend to one.

Choose

\[
L=\frac{(N-1)q}{N^2}.
\]

Then (9) gives

\[
\limsup_{q\to\infty}q^{N-1}C_{\rm opt}(N,N,q)
\le\sigma_{N-1}\left(\frac{eN^2}{N-1}\right)^{N-1}.
\tag{10}
\]

Equations (7) and (10) prove (1).

\hypertarget{1016-additional-bound-for-the-supremum-norm}{%
\section{Additional bound for the supremum
norm}\label{1016-additional-bound-for-the-supremum-norm}}

\hypertarget{1016-statement-and-scope}{%
\subsection{Statement and scope}\label{1016-statement-and-scope}}

Let \(N\geq2\). For a nonempty bounded open convex set
\(\Omega\subset\mathbb R^N\), define its Euclidean inradius and its
Dirichlet Poincaré--Sobolev frequency at \(q=\infty\) by

\[
R_\Omega=\sup\{r>0:\text{there is }x\in\mathbb R^N
\text{ with }B_r(x)\subset\Omega\},
\]

\[
\lambda_{p,\infty}(\Omega)
=\inf_{\substack{f\in C_c^\infty(\Omega;\mathbb R)\\f\ne0}}
\frac{\displaystyle\int_\Omega|\nabla f|^p\,dx}
     {\|f\|_{L^\infty(\Omega)}^p}.
\]

For \(p>N\), the constant requested in this parameter regime of Problem
1016 is

\[
C_{\mathrm{opt}}(N,p,\infty)
=\inf_{\substack{\Omega\subset\mathbb R^N\\
\Omega\text{ nonempty, bounded, open, convex}}}
R_\Omega^{p-N}\lambda_{p,\infty}(\Omega).
\tag{S1}
\]

Scaling gives the equivalent definition with the constraint
\(R_\Omega=1\). The following explicit lower bound holds in the range
\(p\geq N+1\):

\[
\boxed{
C_{\mathrm{opt}}(N,p,\infty)
\geq
2\left(\frac{p-N}{p-1}\right)^{p-1}
\frac{\pi^{(N-1)/2}\Gamma((p-N+1)/2)}{\Gamma(p/2)}.
}
\tag{S2}
\]

For example, (S2) gives the exact comparison

\[
C_{\mathrm{opt}}(2,4,\infty)\geq\frac{8\pi}{27}.
\]

This is a lower bound, not a determination of the optimal constant. No
equality case, minimizing domain, or sharpness assertion is made. It
adds a parameter regime different from the critical-exponent asymptotic
\(p=N\), \(q\to\infty\). The proof below does not cover \(N<p<N+1\), and
no claim of bibliographic priority is made.

\hypertarget{1016-balanced-supporting-planes}{%
\subsection{1. Balanced supporting
planes}\label{1016-balanced-supporting-planes}}

Normalize \(R_\Omega=1\), and translate a center of a maximal inscribed
ball to the origin. Such a center exists because \(\Omega\) is bounded.
There are contact supporting planes with unit outward normals
\(n_1,\ldots,n_m\), and positive weights \(\lambda_1,\ldots,\lambda_m\),
such that

\[
\Omega\subset P:=\bigcap_{i=1}^m\{x:n_i\cdot x<1\},
\qquad
\sum_{i=1}^m\lambda_i=1,
\qquad
\sum_{i=1}^m\lambda_i n_i=0.
\tag{S3}
\]

Here one can take \(m\leq N+1\). To justify the balance condition, let
\(h_\Omega(n)=\sup_{x\in\Omega}n\cdot x\), and consider the compact set
of contact normals satisfying \(h_\Omega(n)=1\). If their convex hull
did not contain zero, strict separation would give a direction in which
moving the ball center increases the distance from all contact planes.
The support function is continuous, so the same strict gain holds in a
neighborhood of the contact normals; away from that neighborhood the
support distances already exceed one by a positive amount. A
sufficiently small displacement would therefore permit a larger
inscribed ball, a contradiction. Carathéodory's theorem gives the finite
collection in (S3); zero weights can be discarded.

Fix a nonzero test function \(f\in C_c^\infty(\Omega;\mathbb R)\).
Choose a point \(x_*\in\Omega\) at which its absolute value attains its
maximum, and put

\[
M=|f(x_*)|=\|f\|_\infty,
\qquad
a_i=1-n_i\cdot x_*>0,
\qquad
v_i=\frac{n_i}{a_i},
\qquad
\rho_i=\lambda_i a_i.
\]

The identities in (S3) imply

\[
\sum_i\lambda_i a_i=1,
\qquad
\sum_i\rho_i=1,
\qquad
\sum_i\rho_i v_i=0.
\tag{S4}
\]

Thus the \(\rho_i\) are probability weights, and the vectors \(v_i\)
have weighted mean zero. In particular, the maximum point \(x_*\) need
not be an incenter.

\hypertarget{1016-a-support-function-estimate}{%
\subsection{2. A support-function
estimate}\label{1016-a-support-function-estimate}}

Write

\[
s=p-N\geq1,
\qquad
h(\omega)=\max_i(v_i\cdot\omega),
\quad\omega\in S^{N-1}.
\]

Equation (S4) implies \(h\geq0\). Along the ray \(x_*+r\omega\), the
exit length from \(P\) is \(1/h(\omega)\), with value \(+\infty\) when
\(h(\omega)=0\).

For a fixed direction \(\omega\), set

\[
A=h(\omega),\qquad B=h(-\omega),\qquad X_i=v_i\cdot\omega.
\]

Then \(-B\leq X_i\leq A\) and \(\sum_i\rho_iX_i=0\). Convexity of
\(x\mapsto|x|^s\) and comparison with its chord over \([-B,A]\) give

\[
\sum_i\rho_i|X_i|^s
\leq\frac{AB^s+BA^s}{A+B}
\leq\frac{A^s+B^s}{2}.
\tag{S5}
\]

The second inequality is equivalent to \((A-B)(A^s-B^s)\geq0\). If
\(A+B=0\), all terms vanish and (S5) holds directly.

Integrating (S5) over the unit sphere and using its symmetry yields

\[
\int_{S^{N-1}}h(\omega)^s\,d\omega
\geq\kappa_{N,s}\sum_i\rho_i|v_i|^s
=\kappa_{N,s}\sum_i\lambda_i a_i^{1-s},
\tag{S6}
\]

where the beta integral gives, for any unit vector \(e\),

\[
\kappa_{N,s}
=\int_{S^{N-1}}|e\cdot\omega|^s\,d\omega
=\frac{2\pi^{(N-1)/2}\Gamma((s+1)/2)}
       {\Gamma((N+s)/2)}.
\tag{S7}
\]

Since \(s\geq1\), the function \(a\mapsto a^{1-s}\) is convex on
\((0,\infty)\). Jensen's inequality and (S4) therefore imply

\[
\sum_i\lambda_i a_i^{1-s}
\geq\left(\sum_i\lambda_i a_i\right)^{1-s}=1.
\]

Consequently

\[
\int_{S^{N-1}}h(\omega)^s\,d\omega\geq\kappa_{N,s}.
\tag{S8}
\]

\hypertarget{1016-radial-energy-from-the-maximum-point}{%
\subsection{3. Radial energy from the maximum
point}\label{1016-radial-energy-from-the-maximum-point}}

Let \(R(\omega)\) be the exit length from \(\Omega\) along
\(x_*+r\omega\). Convexity ensures that this ray stays in the domain for
\(0\leq r<R(\omega)\), and (S3) gives

\[
R(\omega)\leq\frac1{h(\omega)}.
\tag{S9}
\]

The test function vanishes before the exit. The fundamental theorem of
calculus, followed by weighted Hölder, gives

\[
\begin{aligned}
M
&\leq\int_0^{R(\omega)}|\partial_r f(x_*+r\omega)|\,dr\\
&\leq
\left(\int_0^{R(\omega)}|\partial_r f|^p r^{N-1}\,dr\right)^{1/p}
\left(\int_0^{R(\omega)}r^{-(N-1)/(p-1)}\,dr\right)^{(p-1)/p}.
\end{aligned}
\]

The last integral is finite because \(p>N\). Evaluating it and
rearranging give

\[
\int_0^{R(\omega)}|\partial_r f|^p r^{N-1}\,dr
\geq
\left(\frac{p-N}{p-1}\right)^{p-1}
M^pR(\omega)^{N-p}.
\tag{S10}
\]

Since \(N-p=-s\), (S9) implies \(R(\omega)^{N-p}\geq h(\omega)^s\),
including the case \(h(\omega)=0\). Integrate (S10) in direction, use
\(|\partial_r f|\leq|\nabla f|\), and apply (S8). This proves

\[
\frac{\int_\Omega|\nabla f|^p\,dx}{\|f\|_\infty^p}
\geq
\left(\frac{p-N}{p-1}\right)^{p-1}\kappa_{N,p-N}.
\]

Taking the infimum over all tests and all convex domains of inradius
one, and substituting (S7), gives (S2). For a general inradius, scaling
restores the factor \(R_\Omega^{p-N}\) in (S1).

The two convexity steps in Section 2 require \(p-N\geq1\). Moreover, the
radial estimate omits angular derivatives. Equality in an auxiliary
geometric inequality therefore would not establish sharpness of (S2) for
the original shape optimization problem.

\section{Problem source and prior asymptotics}
The sharp-inradius question is posed in Brasco's Open problems~1,
following the superhomogeneous inequality in Proposition~6.3.
For a fixed planar domain, the large-\(q\) coefficient \(8\pi e\)
also appears in Bozzola--Brasco, Corollary~3.7 of the published version.
A fixed-domain limit alone does not justify interchanging the limit with
an infimum over all convex shapes; the uniform Hardy and
Adachi--Tanaka argument above provides the needed uniformity.
\begin{enumerate}
\raggedright
\item L. Brasco, \emph{On principal frequencies and inradius in convex sets},
Bruno Pini Mathematical Analysis Seminar \textbf{9}(1) (2018), 78--101,
Proposition~6.3 and (15), p.~99; Open problems~1, p.~100.
\url{https://doi.org/10.6092/issn.2240-2829/8945}.
\url{https://mathematicalanalysis.unibo.it/article/download/8945/8866/28804}.
\item F. Bozzola and L. Brasco,
\emph{The role of topology and capacity in some bounds for principal frequencies},
The Journal of Geometric Analysis \textbf{34} (2024), Article~299.
Published Section~3.3, Corollary~3.7, p.~23.
\url{https://doi.org/10.1007/s12220-024-01742-2}.
The author manuscript at
\url{https://cvgmt.sns.it/media/doc/paper/6357/bozbra_final.pdf}
numbers the same result Corollary~3.8, pp.~17--18.
\end{enumerate}
\end{document}
